% GENERATED FILE. Edit canonical lessons, then run npm run generate:books.
% canonical-source-hash: fnv1a64:82cff105753da324
% canonical-lessons: JA-C295-hatsumoude, JA-C295-omikuji, JA-C295-omamori, JA-C295-narasu, JA-C295-hibiku

\chapter{New Year at the Shrine}
\label{ch:ja-new-year-at-the-shrine}
\hlchaptermodality{295}

\begin{chapteropening}
\textbf{By the end of this chapter:} \emph{I can say what we did at the shrine at New Year.}

Complete the last lesson of chapter 295: I can say what we did at the shrine at New Year.
\end{chapteropening}

\section[\texorpdfstring{\ja{はつもうで} (hatsumōde)}{はつもうで (hatsumōde)}]{\ja{はつもうで} (hatsumōde) — the New Year shrine visit; the first shrine visit of the year}

\label{lesson:JA-C295-hatsumoude}



\begin{quote}
Before the new one: say the Japanese for to liven up, then the Japanese for a Bon festival dance.
\end{quote}

\subsection*{You'll want to know: \ja{はつもうで}}
\textbf{\ja{はつもうで}} — \emph{hatsumōde} — \textquotedblleft{}the New Year shrine visit; the first shrine visit of the year\textquotedblright{}.

\begin{quote}
\textbf{\ja{しょうがつに} \ja{かぞくで} \ja{はつもうでに} \ja{いきました。}} — \emph{shōgatsu ni kazoku de hatsumōde ni ikimashita} — \textquotedblleft{}At the new year the family made its first shrine visit.\textquotedblright{}
\end{quote}

\begin{cousinweb}
\textbf{\ja{はつもうで}} (\emph{hatsumōde}) is a native Japanese word (wago), inherited rather than borrowed from Chinese, built from \emph{hatsu}, \textquotedblleft{}first\textquotedblright{}, and \emph{mōde}, \textquotedblleft{}a visit to a shrine\textquotedblright{} (JMdict, KANJIDIC2).
\end{cousinweb}

\begin{grammarlens}[title={What you've built}]
One more word for telling what happened.
\end{grammarlens}

\subsection*{Your turn}
\begin{itemize}
\raggedright
  \item \emph{Say it:} \emph{hatsumōde}
  \item \emph{Say it:} \emph{hatsumōde}, once more
  \item \emph{Say it:} \emph{shōgatsu ni kazoku de hatsumōde ni ikimashita} — At the new year the family made its first shrine visit.
  \item \emph{Recall:} say the Japanese for to liven up, then the Japanese for a Bon festival dance, then say \emph{hatsumōde} again
\end{itemize}

\begin{tcolorbox}[breakable,colback=teal!4,colframe=teal!35!black,title={Before you move on}]
What does \ja{はつもうで} mean? (\textquotedblleft{}The New Year shrine visit; the first shrine visit of the year\textquotedblright{}.) Say it once more.
\end{tcolorbox}
\section[\texorpdfstring{\ja{おみくじ} (omikuji)}{おみくじ (omikuji)}]{\ja{おみくじ} (omikuji) — a shrine fortune slip; a paper fortune}

\label{lesson:JA-C295-omikuji}



\begin{quote}
Before the new one: say the Japanese for a Bon festival dance, then the Japanese for the New Year shrine visit.
\end{quote}

\subsection*{You'll want to know: \ja{おみくじ}}
\textbf{\ja{おみくじ}} — \emph{omikuji} — \textquotedblleft{}a shrine fortune slip; a paper fortune\textquotedblright{}.

\begin{quote}
\textbf{\ja{じんじゃで} \ja{おみくじを} \ja{ひきました。}} — \emph{jinja de omikuji o hikimashita} — \textquotedblleft{}I drew a fortune at the shrine.\textquotedblright{}
\end{quote}

\begin{cousinweb}
\textbf{\ja{おみくじ}} (\emph{omikuji}) is a native Japanese word (wago), inherited rather than borrowed from Chinese, built from \emph{o-}, \textquotedblleft{}the polite prefix\textquotedblright{}, \emph{mi-}, \textquotedblleft{}the honorific prefix\textquotedblright{}, and \emph{kuji}, \textquotedblleft{}a lot drawn\textquotedblright{} (JMdict, KANJIDIC2).
\end{cousinweb}

\begin{grammarlens}[title={What you've built}]
One more word for telling what happened.
\end{grammarlens}

\subsection*{Your turn}
\begin{itemize}
\raggedright
  \item \emph{Say it:} \emph{omikuji}
  \item \emph{Say it:} \emph{omikuji}, once more
  \item \emph{Say it:} \emph{jinja de omikuji o hikimashita} — I drew a fortune at the shrine.
  \item \emph{Recall:} say the Japanese for a Bon festival dance, then the Japanese for the New Year shrine visit, then say \emph{omikuji} again
\end{itemize}

\begin{tcolorbox}[breakable,colback=teal!4,colframe=teal!35!black,title={Before you move on}]
What does \ja{おみくじ} mean? (\textquotedblleft{}A shrine fortune slip; a paper fortune\textquotedblright{}.) Say it once more.
\end{tcolorbox}
\section[\texorpdfstring{\ja{おまもり} (omamori)}{おまもり (omamori)}]{\ja{おまもり} (omamori) — a charm; an amulet}

\label{lesson:JA-C295-omamori}



\begin{quote}
Before the new one: say the Japanese for the New Year shrine visit, then the Japanese for a shrine fortune slip.
\end{quote}

\subsection*{You'll want to know: \ja{おまもり}}
\textbf{\ja{おまもり}} — \emph{omamori} — \textquotedblleft{}a charm; an amulet\textquotedblright{}.

\begin{quote}
\textbf{\ja{ははが} \ja{じんじゃで} \ja{おまもりを} \ja{かいました。}} — \emph{haha ga jinja de omamori o kaimashita} — \textquotedblleft{}My mother bought a charm at the shrine.\textquotedblright{}
\end{quote}

\begin{cousinweb}
\textbf{\ja{おまもり}} (\emph{omamori}) is a native Japanese word (wago), inherited rather than borrowed from Chinese, built from \emph{o-}, \textquotedblleft{}the polite prefix\textquotedblright{}, and \emph{mamori}, \textquotedblleft{}protection\textquotedblright{} (from \emph{mamoru}) (JMdict, KANJIDIC2).
\end{cousinweb}

\begin{grammarlens}[title={What you've built}]
One more word for telling what happened.
\end{grammarlens}

\subsection*{Your turn}
\begin{itemize}
\raggedright
  \item \emph{Say it:} \emph{omamori}
  \item \emph{Say it:} \emph{omamori}, once more
  \item \emph{Say it:} \emph{haha ga jinja de omamori o kaimashita} — My mother bought a charm at the shrine.
  \item \emph{Recall:} say the Japanese for the New Year shrine visit, then the Japanese for a shrine fortune slip, then say \emph{omamori} again
\end{itemize}

\begin{tcolorbox}[breakable,colback=teal!4,colframe=teal!35!black,title={Before you move on}]
What does \ja{おまもり} mean? (\textquotedblleft{}A charm; an amulet\textquotedblright{}.) Say it once more.
\end{tcolorbox}
\section[\texorpdfstring{\ja{ならす} (narasu)}{ならす (narasu)}]{\ja{ならす} (narasu) — to ring; to sound}

\label{lesson:JA-C295-narasu}



\begin{quote}
Before the new one: say the Japanese for a shrine fortune slip, then the Japanese for a charm.
\end{quote}

\subsection*{You'll want to know: \ja{ならす}}
\textbf{\ja{ならす}} — \emph{narasu} — \textquotedblleft{}to ring; to sound\textquotedblright{}.

\begin{quote}
\textbf{\ja{てらで} \ja{かねを} \ja{ならしました。}} — \emph{tera de kane o narashimashita} — \textquotedblleft{}We rang the bell at the temple.\textquotedblright{}
\end{quote}

Polite \textbf{\ja{ならします}}, \emph{narashimasu}; a godan verb. Past (polite): \textbf{\ja{おとこが} \ja{かねを} \ja{ならしました。}}, \emph{otoko ga kane o narashimashita}.

\begin{cousinweb}
\textbf{\ja{ならす}} (\emph{narasu}) is a native Japanese word (wago), inherited rather than borrowed from Chinese; its kanji was attached later as a kun'yomi (JMdict).
\end{cousinweb}

\begin{grammarlens}[title={What you've built}]
One more word for telling what happened.
\end{grammarlens}

\subsection*{Your turn}
\begin{itemize}
\raggedright
  \item \emph{Say it:} \emph{narasu}
  \item \emph{Say it:} \emph{narasu}, once more
  \item \emph{Say it:} \emph{tera de kane o narashimashita} — We rang the bell at the temple.
  \item \emph{Recall:} say the Japanese for a shrine fortune slip, then the Japanese for a charm, then say \emph{narasu} again
\end{itemize}

\begin{tcolorbox}[breakable,colback=teal!4,colframe=teal!35!black,title={Before you move on}]
What does \ja{ならす} mean? (\textquotedblleft{}To ring; to sound\textquotedblright{}.) Say it once more.
\end{tcolorbox}
\section[\texorpdfstring{\ja{ひびく} (hibiku)}{ひびく (hibiku)}]{\ja{ひびく} (hibiku) — to ring out; to echo}

\label{lesson:JA-C295-hibiku}



\begin{quote}
Before the new one: say the Japanese for a charm, then the Japanese for to ring.
\end{quote}

\subsection*{You'll want to know: \ja{ひびく}}
\textbf{\ja{ひびく}} — \emph{hibiku} — \textquotedblleft{}to ring out; to echo\textquotedblright{}.

\begin{quote}
\textbf{\ja{ゆうべ、てらの} \ja{かねが} \ja{ひびいていました。}} — \emph{yūbe, tera no kane ga hibiite imashita} — \textquotedblleft{}Last night the temple bell was ringing out.\textquotedblright{}
\end{quote}

Polite \textbf{\ja{ひびきます}}, \emph{hibikimasu}; a godan verb. Past (polite): \textbf{\ja{かねの} \ja{おとが} \ja{ひびきました。}}, \emph{kane no oto ga hibikimashita}.

\begin{cousinweb}
\textbf{\ja{ひびく}} (\emph{hibiku}) is a native Japanese word (wago), inherited rather than borrowed from Chinese; its kanji was attached later as a kun'yomi (JMdict).
\end{cousinweb}

\begin{grammarlens}[title={What you've built}]
One more word for telling what happened.
\end{grammarlens}

\subsection*{Your turn}
\begin{itemize}
\raggedright
  \item \emph{Say it:} \emph{hibiku}
  \item \emph{Say it:} \emph{hibiku}, once more
  \item \emph{Say it:} \emph{yūbe, tera no kane ga hibiite imashita} — Last night the temple bell was ringing out.
  \item \emph{Recall:} say the Japanese for a charm, then the Japanese for to ring, then say \emph{hibiku} again
\end{itemize}

\begin{tcolorbox}[breakable,colback=teal!4,colframe=teal!35!black,title={Before you move on}]
What does \ja{ひびく} mean? (\textquotedblleft{}To ring out; to echo\textquotedblright{}.) Say it once more.
\end{tcolorbox}
